\documentclass[11pt,letterpaper]{article}

\usepackage[margin=1in,headheight=14pt]{geometry}
\usepackage{amsmath,amssymb,amsthm,mathtools}
\usepackage{graphicx}
\usepackage{xcolor}
\usepackage{enumitem}
\usepackage{microtype}
\usepackage{fancyhdr}
\usepackage[colorlinks=true,linkcolor=blue,citecolor=blue,urlcolor=blue]{hyperref}

% Personal notation and theorem environments.
%!TEX root = EM_singular_models.tex

% PDF margin etc settings
%\setlength{\textwidth}{\paperwidth}
%\addtolength{\textwidth}{-6cm}
%\setlength{\textheight}{\paperheight}
%\addtolength{\textheight}{-4cm}
%\addtolength{\textheight}{-1.1\headheight}
%\addtolength{\textheight}{-\headsep}
%\addtolength{\textheight}{-\footskip}
%\setlength{\oddsidemargin}{0.5cm}
%\setlength{\evensidemargin}{0.5cm}

%%%%%%%%%%%%%%%%%%%%%%%%%%%%%%%%

% MACROS HERE

%%%%%%%%%%%%%%%%%%%%%%%%%%%%%%%%

\newcommand{\fixedpointtheta}{\widehat\theta_\obs}

\newcommand{\cover}{N}
\newcommand{\kinit}{\nu}
\newcommand{\cdelta}{c_\threshold}
\newcommand{\dndelta}{\omega}
\newcommand{\kndelta}{\Upsilon}
\newcommand{\kndeltanew}{\Psi}
\newcommand{\radii}{\mathcal{R}}
\newcommand{\event}{\mathcal{E}}
\newcommand{\kevent}{\mathcal{D}}


% Observations, dimension etc.
\newcommand{\dims}{\ensuremath{d}}
\newcommand{\samples}{\ensuremath{n}}
\newcommand{\obs}{\samples}
\newcommand{\real}{\ensuremath{\mathbb{R}}}
\newcommand{\complex}{\ensuremath{\mathbb{C}}}
\newcommand{\realdim}{\ensuremath{\real^\dims}}

\newcommand{\smallthreshold}{\epsilon}

% some mathcal notations
\newcommand{\order}[1]{\ensuremath{\mathcal{O}\parenth{#1}}}

% Basic statistics notation
\newcommand{\xstar}{\ensuremath{x^*}}
\newcommand{\xbar}{\ensuremath{\bar{x}}}
% True parameter
\newcommand{\thetastar}{\ensuremath{{\theta^*}}}
% Estimate one
\newcommand{\thetahat}{\ensuremath{\widehat{\theta}}}
% Estimate two
\newcommand{\thetatil}{\ensuremath{\widetilde{\theta}}}
\newcommand{\thetabar}{\ensuremath{\bar{\theta}}}
\newcommand{\yvec}{\ensuremath{y}}
\newcommand{\Xmat}{\ensuremath{X}}

% Distributions
\newcommand{\distribution}{P}
\newcommand{\density}{f}
\newcommand{\gaussiandensity}{\phi}
\newcommand{\NORMAL}{\ensuremath{\mathcal{N}}}
\newcommand{\BER}{\ensuremath{\mbox{Ber}}}

% Spaces
\newcommand{\Xspace}{\ensuremath{\mathcal{X}}}
\newcommand{\Yspace}{\ensuremath{\mathcal{Y}}}
\newcommand{\Zspace}{\ensuremath{\mathcal{Z}}}


% Brackets Size
\newcommand{\brackets}[1]{\left[ #1 \right]}
\newcommand{\parenth}[1]{\left( #1 \right)}
\newcommand{\bigparenth}[1]{\big( #1 \big)}
\newcommand{\biggparenth}[1]{\bigg( #1 \bigg)}
\newcommand{\braces}[1]{\left\{ #1 \right \}}
\newcommand{\abss}[1]{\left| #1 \right |}
\newcommand{\angles}[1]{\left\langle #1 \right \rangle}
\newcommand{\tp}{^\top}

% Generic vectors and scalars
\newcommand{\myvec}{u} % for generic vector
\newcommand{\myvectwo}{v} % for generic vector
\newcommand{\mymat}{B}
\newcommand{\set}{{A}}
\newcommand{\scalar}{\alpha}
\newcommand{\tvscalar}{\epsilon}
\newcommand{\tailboundscalar}{t}
\newcommand{\term}{U}
\newcommand{\myVarone}{\alpha_X}
\newcommand{\myVartwo}{\alpha_Y}
\newcommand{\myCovmat}{\Sigma}
\newcommand{\xcord}[1]{x^{(#1)}}
\newcommand{\xcordpop}[1]{X^{(#1)}}
\newcommand{\ycordpop}[1]{Y^{(#1)}}
\newcommand{\myVarthree}{\ensuremath{U}}
\newcommand{\myVarfour}{\ensuremath{V}}
\newcommand{\myfuntwo}{\ensuremath{g}}
\newcommand{\numsteps}{t}
\newcommand{\totalnumsteps}{T}
\newcommand{\seqalpha}[1]{\alpha_{#1}}
\newcommand{\seqnu}[1]{\nu^{(#1)}}
\newcommand{\imax}{\ell_\smallthreshold}
\newcommand{\basis}{e}
\newcommand{\radem}{\varepsilon}
\newcommand{\Tone}{T\ensuremath{_{0}}}
\newcommand{\Ttwo}{T\ensuremath{_{2}}}
\newcommand{\contractPar}[1]{\ensuremath{\gamma\parenth{#1}}}

\newcommand{\simiid}{\overset{\mathrm{i.i.d.}}{\sim}}
\newcommand{\eqd}{\overset{d}{=}}


% EM updates
\newcommand{\hidden}{Z}
\newcommand{\qfun}[2]{Q(#1\vert#2)}
\newcommand{\variable}{\theta}
\newcommand{\weights}{w}
\newcommand{\myfun}{q}
\newcommand{\gradf}{\nabla \myfun}
\newcommand{\hessf}{\nabla^2 \myfun}
\newcommand{\iter}{t}
\newcommand{\mupdate}{M}
\newcommand{\mupdategeneral}{\ensuremath{M^{GN}}}




% EM contractions
\newcommand{\contractgene}{\ensuremath{\gamma^{GN}}}
\newcommand{\updateMat}{\ensuremath{\Gamma_{\theta}}}
\newcommand{\mymatTheta}{\ensuremath{B_{\theta}}}
\newcommand{\sech}{\text{sech}}

% Nhat's macros 
\newcommand{\Rspace}{\ensuremath{\mathbb{R}}}
\newcommand{\Ecal}{\ensuremath{\mathcal{E}}}
\newcommand{\Ncal}{\ensuremath{\mathcal{N}}}
\newcommand{\gradient}{\ensuremath{\nabla}}
\newcommand{\smallweight}{\ensuremath{\epsilon^{sw}}}
\newcommand{\smallnorm}{\ensuremath{\rho^{sw}}}
\newcommand{\ball}{\ensuremath{\mathbb{B}}}
\newcommand{\sphere}{\ensuremath{\mathbb{S}}}

\newcommand{\diam}{\ensuremath{\text{diam}}}
\newcommand{\upper}{\ensuremath{\bar{C}}}
\newcommand{\sample}{\ensuremath{\xi}}
\newcommand{\weight}{\ensuremath{\pi}}
\newcommand{\contracasym}{\ensuremath{\rho}}
\newcommand{\Fishermatrix}{\ensuremath{\mathcal{I}}}
\newcommand{\barcontracasym}{\ensuremath{\bar{\rho}}}


\newcommand{\TrueDensity}{\ensuremath{f}}
\newcommand{\FitDensity}{\ensuremath{f}_{\theta}}
\newcommand{\FitDensityfree}{\ensuremath{f}_{\pi, \theta}}
\newcommand{\FitDensitygen}{\ensuremath{f}_{\pi, \theta_{1}, \theta_{2}}}
\newcommand{\FitDensitymore}{\ensuremath{f}_{G}}
\newcommand{\normden}{\ensuremath{\phi}}
\newcommand{\mean}{\ensuremath{\theta}}
\newcommand{\sd}{\ensuremath{\sigma}}
\newcommand{\parFOS}{\ensuremath{\gamma}}
\newcommand{\funQ}{\ensuremath{Q}}
\newcommand{\updateM}{\ensuremath{M}}
\newcommand{\weightFun}{\ensuremath{w}}
\newcommand{\mydefn}{\ensuremath{:=}}
\newcommand{\poincare}{\ensuremath{\tilde{C}}}
\newcommand{\paraspace}{\ensuremath{\Theta}}
\newcommand{\loglihood}{\ensuremath{F}}
\newcommand{\prior}{\ensuremath{\pi}}
\newcommand{\posterior}{\ensuremath{\Pi}}
\newcommand{\boundcons}{\ensuremath{B}}
\newcommand{\boundconstot}{\ensuremath{H}}
\newcommand{\noise}{\ensuremath{\varepsilon}}
\newcommand{\loglihoodgaus}{\ensuremath{F^{G}}}
\newcommand{\loglihoodind}{\ensuremath{F^{I}}}
\newcommand{\loglihoodlogit}{\ensuremath{F^{R}}}
\newcommand{\noiseone}{\ensuremath{\varepsilon_{1}}}
\newcommand{\noisetwo}{\ensuremath{\varepsilon_{2}}}

%%%SDE_Bayesian_inf
\newcommand{\weakcon}{\ensuremath{\psi}}
\newcommand{\perturb}{\ensuremath{\zeta}}
\newcommand{\inver}{\ensuremath{\xi}}

% Universal constants
\newcommand{\unicon}{c}
\newcommand{\uniconnew}{c'}
\newcommand{\uniconone}{c_1}
\newcommand{\unicontwo}{c_2}

% some parameters that you may define
\newcommand{\scparam}{m} % strong  convexity parameter
\newcommand{\smoothness}{L} % smoothness parameter
\newcommand{\step}{h}


%%%%%%%%% Distributions and Random variables %%%%%%%%%%%
\newcommand{\rvg}{\xi} % to denote the random variable g
\newcommand{\threshold}{\delta}


%%%%%%%%% Basic Terms like defn, etal, tmix, polylog %%%%%%%%%%%
\newcommand{\defn}{:=}
\newcommand{\rdefn}{=:}
\newcommand{\etal}{{et al.}}
\newcommand{\polylog}{\text{poly-log}}

% Some vector/matrix norms
\newcommand{\matnorm}[3]{|\!|\!| #1 | \! | \!|_{{#2}, {#3}}}
\newcommand{\matsnorm}[2]{|\!|\!| #1 | \! | \!|_{{#2}}}
\newcommand{\vecnorm}[2]{\left\| #1\right\|_{#2}}
\newcommand{\enorm}[1]{\vecnorm{#1}{2}} % euclidean norm
\newcommand{\nucnorm}[1]{\ensuremath{\matsnorm{#1}{\footnotesize{\mbox{nuc}}}}}
\newcommand{\fronorm}[1]{\ensuremath{\matsnorm{#1}{\footnotesize{\mbox{fro}}}}}
\newcommand{\opnorm}[1]{\ensuremath{\matsnorm{#1}{\tiny{\mbox{op}}}}}
%%% EM paper
\newcommand{\constrainnorm}[1]{\vecnorm{#1}{[k]}} % euclidean norm

\newcommand{\bigmatsnorm}[2]{{\Big|\!\Big|\!\Big| #1
  \Big|\!\Big|\!\Big|}_{#2}}

% Inner product
\newcommand{\inprod}[2]{\ensuremath{\langle #1 , \, #2 \rangle}}

% Kullback-Leibler
\newcommand{\kull}[2]{\ensuremath{D_{\text{KL}}(#1\; \| \; #2)}}

% Probability
\newcommand{\Exs}{\ensuremath{{\mathbb{E}}}}
\newcommand{\Prob}{\ensuremath{{\mathbb{P}}}}

%Eigenvector / eigenvalue related notation
\newcommand{\eig}[1]{\ensuremath{\lambda_{#1}}}
\newcommand{\eigmax}{\ensuremath{\eig{\max}}}
\newcommand{\eigmin}{\ensuremath{\eig{\min}}}

% \DeclareMathOperator{\det}{det}
\DeclareMathOperator{\modd}{mod}
\DeclareMathOperator{\diag}{diag}
\DeclareMathOperator{\var}{var}
\DeclareMathOperator{\cov}{cov}
\DeclareMathOperator{\trace}{trace}
\DeclareMathOperator{\abs}{abs}
\DeclareMathOperator{\floor}{floor}
\DeclareMathOperator{\vol}{vol}
\DeclareMathOperator{\child}{child}
\DeclareMathOperator{\parent}{parent}
\DeclareMathOperator{\sign}{sign}
\DeclareMathOperator{\rank}{rank}
\DeclareMathOperator{\toeplitz}{toeplitz}




\newtheoremstyle{named}{}{}{\itshape}{}{\bfseries}{.}{.5em}{\thmnote{#3's }#1}
\theoremstyle{named}
\newtheorem*{namedtheorem}{Lemma}

%%%%%%%
\theoremstyle{plain}

% {Theorem, Proposition, Lemma, Corollary} numbered sequentially
% throughout the paper
\newtheorem{theorem}{Theorem}
\newtheorem{proposition}{Proposition}
\newtheorem{lemma}{Lemma}
\newtheorem{claim}{Claim}
\newtheorem{corollary}{Corollary}
\newtheorem{definition}{Definition}
\newtheorem{conjecture}{Conjecture}
\newtheorem{remark}{Remark}
%%%%%%%%%%%%%%%%%%%%%%%%%%%%%%%%%%%%%%%%%%%%%%%%%%%%%%%%%%%%%%%%%%%%%%%
% WIDEBAR COMMAND
\newlength{\widebarargwidth}
\newlength{\widebarargheight}
\newlength{\widebarargdepth}
\DeclareRobustCommand{\widebar}[1]{%
  \settowidth{\widebarargwidth}{\ensuremath{#1}}%
  \settoheight{\widebarargheight}{\ensuremath{#1}}%
  \settodepth{\widebarargdepth}{\ensuremath{#1}}%
  \addtolength{\widebarargwidth}{-0.3\widebarargheight}%
  \addtolength{\widebarargwidth}{-0.3\widebarargdepth}%
  \makebox[0pt][l]{\hspace{0.3\widebarargheight}%
    \hspace{0.3\widebarargdepth}%
    \addtolength{\widebarargheight}{0.3ex}%
    \rule[\widebarargheight]{0.95\widebarargwidth}{0.1ex}}%
  {#1}}



%%% New version of \caption puts things in smaller type, single-spaced
%%% and indents them to set them off more from the text.
\makeatletter
\long\def\@makecaption#1#2{
        \vskip 0.8ex
        \setbox\@tempboxa\hbox{\small {\bf #1:} #2}
        \parindent 1.5em  %% How can we use the global value of this???
        \dimen0=\hsize
        \advance\dimen0 by -3em
        \ifdim \wd\@tempboxa >\dimen0
                \hbox to \hsize{
                        \parindent 0em
                        \hfil
                        \parbox{\dimen0}{\def\baselinestretch{0.96}\small
                                {\bf #1.} #2
                                %%\unhbox\@tempboxa
                                }
                        \hfil}
        \else \hbox to \hsize{\hfil \box\@tempboxa \hfil}
        \fi
        }
\makeatother



%% COMMENTING commands

\long\def\comment#1{}
\definecolor{battleshipgrey}{rgb}{0.52, 0.52, 0.51}
\definecolor{darkgray}{rgb}{0.66, 0.66, 0.66}
\definecolor{darkgreen}{rgb}{0.0, 0.2, 0.13}
\definecolor{darkspringgreen}{rgb}{0.09, 0.45, 0.27}
\definecolor{dukeblue}{rgb}{0.0, 0.0, 0.61}
\definecolor{olivedrab7}{rgb}{0.24, 0.2, 0.12}
\definecolor{darkblue}{rgb}{0.0, 0.0, 0.55}
\definecolor{darkscarlet}{rgb}{0.34, 0.01, 0.1}
\definecolor{candyapplered}{rgb}{1.0, 0.03, 0.0}
\definecolor{ao(english)}{rgb}{0.0, 0.5, 0.0}
\definecolor{applegreen}{rgb}{0.55, 0.71, 0.0}

\newcommand{\red}[1]{\textcolor{red}{#1}}
\newcommand{\mjwcomment}[1]{{\bf{{\red{{MJW --- #1}}}}}}
\newcommand{\mwlcomment}[1]{{\bf{{\color{orange}{{Wenlong --- #1}}}}}}

\newcommand{\blue}[1]{\textcolor{blue}{#1}}
\newcommand{\green}[1]{\textcolor{green}{#1}}
\newcommand{\highlight}[1]{\textcolor{applegreen}{#1}}

% comment lines
\newcommand{\genecomment}[1]{{\bf{{\blue{{TEAM --- #1}}}}}}
\newcommand{\rrdcomment}[1]{{\bf{{\blue{{RRD --- #1}}}}}}
\newcommand{\nhcomment}[1]{{\bf{{\blue{{NH --- #1}}}}}}
\newcommand{\kkcomment}[1]{{\bf{{\darkgreen{{KK --- #1}}}}}}


\newcommand{\widgraph}[2]{\includegraphics[keepaspectratio,width=#1]{#2}}


\newcommand{\coursenumber}{STA3000F}
\newcommand{\coursetitle}{Advanced Theory of Statistics}
\newcommand{\courseterm}{Fall 2026}
\newcounter{lecture}
\setcounter{lecture}{7}
\newcommand{\lecturenumber}{\arabic{lecture}}
\newcommand{\lecturedate}{October 23, 2026}
\newcommand{\lecturer}{Wenlong Mou}

\pagestyle{fancy}
\fancyhf{}
\fancyhead[L]{\coursenumber: \coursetitle}
\fancyhead[R]{Lecture \lecturenumber}
\fancyfoot[C]{\lecturenumber--\arabic{page}}
\renewcommand{\headrulewidth}{0.4pt}
\renewcommand{\footrulewidth}{0pt}

\usepackage{bm}

\renewcommand{\thepage}{\lecturenumber--\arabic{page}}

\newtheorem{example}{Example}
\numberwithin{theorem}{lecture}
\numberwithin{proposition}{lecture}
\numberwithin{lemma}{lecture}
\numberwithin{claim}{lecture}
\numberwithin{corollary}{lecture}
\numberwithin{definition}{lecture}
\numberwithin{conjecture}{lecture}
\numberwithin{remark}{lecture}
\numberwithin{example}{lecture}
\numberwithin{equation}{lecture}
\counterwithin*{section}{lecture}
\renewcommand{\theHsection}{\arabic{lecture}.\arabic{section}}

\newcommand{\makelecturenotesheader}{%
  \noindent
  \textbf{\coursenumber: \coursetitle}%
  \hfill
  \textbf{\courseterm}

  \begin{center}
    {\Large\bfseries Lecture \lecturenumber: \lecturedate}
  \end{center}

  \noindent
  \textbf{Lecturer:} \lecturer
  \par\medskip
  \hrule
  \bigskip
}

\begin{document}

\makelecturenotesheader
\section{General empirical process bound continued}
Example of the general theorem: consider $\Theta \subseteq \ball (0, R)$ in $d$-dimensional Euclidean space, and consider the function class parametrized by $\theta \in \Theta$:
\begin{align*}
  \mathcal{F} = \{f_\theta: \theta \in \Theta\}
\end{align*}
(e.g. generalized linear model)

We assume that $f_\theta$ is Lipschitz in $\theta$ for every $x$, i.e., there exists $M(x)$ such that
\begin{align*}
  |f_\theta (x) - f_{\theta'} (x)| \leq M(x) \|\theta - \theta'\|_2, \quad \forall \theta, \theta' \in \Theta.
\end{align*}
(We allow the Lipschitz constant to depend on $x$.) We also assume that $\Exs [M^2 (X)] < + \infty$.

In such a case, we can apply the main theorem to bound $\Exs \big[\sup_{f \in \mathcal{F}} |(P_n - P) f| \big]$. To start with, we need to construct an envelope function for the function class $\mathcal{F}$. We can take
\begin{align*}
  G (x) = |f_{\theta_0} (x)| + 2 R M(x),
\end{align*}
for some fixed $\theta_0 \in \Theta$. Then we have
\begin{align*}
  |f_\theta (x)| \leq G (x), \quad \forall \theta \in \Theta.
\end{align*}
Now we need to bound the covering number of $\mathcal{F}$ with respect to the $\mathbb{L}^2 (Q)$ metric, for any distribution $Q$. First, we note that
\begin{align*}
  N ( \mathcal{F}, \mathbb{L}^2 (Q), \varepsilon \|G\|_{\mathbb{L}^2 (Q)}) & \leq N (\mathcal{F}, \mathbb{L}^2 (Q), \varepsilon \|M\|_{\mathbb{L}^2 (Q)} R).
\end{align*}
Now we need to bound the right-hand-side by direct construction of a covering set. Let $(\theta_i)_{i = 1}^N$ be an $\delta$-covering of $\Theta$ with respect to the $\ell_2$ norm. For any $\theta \in \Theta$, there exists some $\theta_i$ such that $\|\theta - \theta_i\|_2 \leq \delta$. Then we have
\begin{align*}
  \vecnorm{f_\theta - f_{\theta_i}}{\mathbb{L}^2 (Q)}^2 = \int | f_\theta (x) - f_{\theta_i} (x)|^2 d Q(x) \leq \int M^2 (x) \|\theta - \theta_i\|_2^2 d Q(x) \leq \delta^2 \|M\|_{\mathbb{L}^2 (Q)}^2.
\end{align*}
That means $(f_{\theta_i})_{i = 1}^N$ is a $\delta \|M\|_{\mathbb{L}^2 (Q)}$-covering of $\mathcal{F}$ with respect to the $\mathbb{L}^2 (Q)$ metric. Therefore, by taking $\delta = \varepsilon R$, we have
\begin{align*}
  N (\mathcal{F}, \vecnorm{\cdot}{\mathbb{L}^2 (Q)}, \varepsilon R \vecnorm{M}{\mathbb{L}^2 (Q)} ) \leq N (\Theta, \|\cdot\|_2, \varepsilon R) \leq N (\ball_2^d (R), \|\cdot\|_2, \varepsilon R) \leq \Big( 1 + \frac{2}{\varepsilon} \Big)^d.
\end{align*}
Substituting this bound into the main theorem, we have
\begin{align*}
  \Exs \Big[\sup_{f \in \mathcal{F}} |(P_n - P) f| \Big] &\leq c \sqrt{\frac{\Exs [G^2 (X)]}{n}}\int_0^{1} \sqrt{\log \big(2\sup_{Q} N (\mathcal{F}; \mathbb{L}^2 (Q), \vecnorm{G}{\mathbb{L}^2 (Q)} \epsilon)\big)} d \epsilon\\
  &\leq c' \sqrt{\frac{\Exs [G^2 (X)]}{n}}\int_0^{1} \sqrt{d\log (1/\epsilon)} d \epsilon\\
  &\leq c'' \cdot \big( \vecnorm{f_{\theta_0}}{\mathbb{L}^2 (\Prob)} + R \vecnorm{M}{\mathbb{L}^2 (\Prob)}  \big)  \sqrt{\frac{d}{n}}.
\end{align*}
This applies to a broad class of Lipschitz models. On the other hand, some function classes may not be continuously parametrized by a finite-dimensional parameter $\theta$.

e.g. consider the function class
\begin{align*}
  \mathcal{F} = \Big\{ \bm{1}_{x \in [\theta_1, \theta_2]}: \theta_1, \theta_2 \in \mathbb{R} \Big\},
\end{align*}
is a very simple function class parametrized by two parameters $\theta_1, \theta_2$, but not continuously. The above argument does not apply. On the other hand, covering/packing number is a robust notion --- as long as the function class is ``simple'', we can still get a good bound.

\section{Combinatorial notion of complexity: VC dimension}
First, let us consider binary functions mapping from $\mathcal{X}$ to $\{0, 1\}$ (we will extend back to real-valued functions later). Consider a class of functions $\mathcal{F}$.

\begin{definition}[shattering]
  A set of points $\{x_1, \ldots, x_n\} \subset \mathcal{X}$ is said to be shattered by $\mathcal{F}$ if for every subset $S \subset \{1, \ldots, n\}$, there exists a function $f \in \mathcal{F}$ such that
  \begin{align*}
    f (x_i) = \begin{cases}
      1, & i \in S\\
      0, & i \notin S
    \end{cases}, \quad i = 1, \ldots, n.
  \end{align*}
\end{definition}
Intuitively, a function class $\mathcal{F}$ needs to be ``rich enough'' to realize all possible labelings of the points in the set $\{x_1, \ldots, x_n\}$, in order to shatter it. From machine learning perspective, you can also think of shattering as a classifier being able to fit any possible labeling of the points in the set $\{x_1, \ldots, x_n\}$.

Based on the notion of shattering, we can define the Vapnik-Chervonenkis (VC) dimension of a function class $\mathcal{F}$.
\begin{definition}[VC dimension]
  The VC dimension of a function class $\mathcal{F}$ is defined as the largest integer $d$ such that there exists a set of $d$ points in $\mathcal{X}$ that can be shattered by $\mathcal{F}$. If no such largest integer exists, we say that the VC dimension is infinite.
\end{definition}
Intuitively, this is the largest number of data points that the binary classifier class can fit perfectly for any possible labeling.

\begin{example}\upshape Indicator of line segments in $\mathbb{R}$: $\mathcal{F} = \{\bm{1}_{x \in [\theta_1, \theta_2]}: \theta_1, \theta_2 \in \mathbb{R}\}$. Consider any 3 points $x_1 < x_2 < x_3$. For the labeling $(1, 0, 1)$, there is no line segment that can realize this labeling. Therefore, the VC dimension of this function class is at most 2. On the other hand, for any 2 points $x_1 < x_2$, we can realize all possible labelings $(0, 0), (0, 1), (1, 0), (1, 1)$ by choosing appropriate line segments. Therefore, the VC dimension of this function class is exactly 2.

More generally, you can consider union of $K$ disjoint intervals in $\mathbb{R}$:
\begin{align*}
  \mathcal{F} = \Big\{\bm{1}_{x \in \bigcup_{k = 1}^K [\theta_{2k - 1}, \theta_{2k}]}: \theta_1 < \theta_2 < \cdots < \theta_{2K} \in \mathbb{R}\Big\}.
\end{align*}
By putting $2K + 1$ points on real line with alternating labeling $(1, 0, 1, 0, \ldots, 0, 1)$, we can see that the VC dimension of this function class is at most $2K$.
\end{example}

\subsection{Structural results on VC dimension}
Idea: If a function class has finite VC dimension, then it is ``simple'' in some sense, and we can bound its covering number.

Key result:
\begin{theorem}[Sauer's lemma]
  Let $\mathcal{F}$ be a class of binary functions with VC dimension $d$. Then for any $n \geq d$, we have
  \begin{align*}
    \Big| \Big\{ (f (x_1), \ldots, f (x_n)) : f \in \mathcal{F} \Big\} \Big| \leq \underbrace{\binom{n}{0} + \binom{n}{1} + \cdots + \binom{n}{d}}_{\Phi_d (n)} \leq \Big( \frac{e n}{d} \Big)^d.
  \end{align*}
\end{theorem}
Remark: consider the cardinality of all possible labelings of $n$ points by a function class $\mathcal{F}$. 
\begin{itemize}
  \item When $n \leq d$, all the $2^n$ possible labelings can be realized by $\mathcal{F}$.
  \item The definition of VC dimension only tells us that for $n = d + 1$, there exists at least one labeling that cannot be realized by $\mathcal{F}$.
  \item Sauer's lemma tells us that once we pass the threshold $n > d$, the number of realizable labelings is bounded by a polynomial in $n$ (with degree $d$), which is much smaller than the total number of possible labelings $2^n$.
\end{itemize}
Proof idea: induction on $(n + d)$. The result is trivial when $n = d$ or when $d = 0$.

By induction hypothesis, let us assume that the conclusion holds for $(n - 1, d - 1)$ and $(n - 1, d)$. Consider the set of $n$ points $\{x_1, \ldots, x_n\}$. We separate the labeling of the last point $x_n$ from the first $n - 1$ points.

Consider any label sequence $b_1, \ldots, b_{n - 1} \in \{0, 1\}^{n - 1}$, there are 3 possibilities:
\begin{enumerate}
\item There is no function $f \in \mathcal{F}$ that realizes the labeling $(b_1, \ldots, b_{n - 1})$ on the first $n - 1$ points.
\item Both labelings $(b_1, \ldots, b_{n - 1}, 0)$ and $(b_1, \ldots, b_{n - 1}, 1)$ can be realized by some functions in $\mathcal{F}$.
\item Only one of the labelings $(b_1, \ldots, b_{n - 1}, 0)$ and $(b_1, \ldots, b_{n - 1}, 1)$ can be realized by some function in $\mathcal{F}$.
\end{enumerate}
The idea is to form a decomposition of the set $\{ (f (x_1), \ldots, f (x_n)) : f \in \mathcal{F} \}$ based on the labeling of the first $n - 1$ points. The first case does not contribute to the cardinality. For the second and third cases, we can decompose the class into $\mathcal{F}_1$ and $\mathcal{F}_2$ according to the following rule (here we abuse the notation a bit, and think of $\mathcal{F}_1$ and $\mathcal{F}_2$ are defined on the domain $\{x_1, \ldots, x_{n}\}$, and equivalently, they are just labeling sequences):
\begin{itemize}
\item In case 2, suppose $f_1, f_2 \in \mathcal{F}$ realize the labelings $(b_1, \ldots, b_{n - 1}, 0)$ and $(b_1, \ldots, b_{n - 1}, 1)$ respectively. We add $f_1$ to $\mathcal{F}_1$ and $f_2$ to $\mathcal{F}_2$. 
\item In case 3, there is only one function $f \in \mathcal{F}$ that realizes the labeling $(b_1, \ldots, b_{n - 1}, b_n)$ for some $b_n \in \{0, 1\}$. We add $f$ to $\mathcal{F}_1$.
\end{itemize}
By this construction, we have
\begin{align*}
  |\mathcal{F} (x_1, \ldots, x_n)| = |\mathcal{F}_1 | + |\mathcal{F}_2 |.
\end{align*}
For the set $\mathcal{F}_1$, we note that this set is formed by mapping all the achievable labelings of the first $n - 1$ points to an element in $\mathcal{F}_1$. So we have
\begin{align*}
  |\mathcal{F}_1| = |\mathcal{F} (x_1, \ldots, x_{n - 1})| \leq \Phi_d (n - 1),
\end{align*}
by induction hypothesis.

For the set $\mathcal{F}_2$, we claim that
\begin{align}
  \text{VCdim} (\mathcal{F}_2) \leq d - 1.\label{eq:lec6-vc-dim-f2}
\end{align}
To prove this claim, we argue by contradiction. Suppose that there exists $(x_{k_1}, x_{k_2}, \cdots, x_{k_d})$ that can be shattered by $\mathcal{F}_2$. Because $\mathcal{F}_2$ is formed by sequences that allow free labeling of the last point $x_n$, we can add this free labeling to the set of $d$ points. This means that $(x_{k_1}, x_{k_2}, \cdots, x_{k_d}, x_n)$ can be shattered by $\mathcal{F}$. This contradicts the assumption that $\mathcal{F}$ has VC dimension $d$. Therefore, we have proved the claim~\eqref{eq:lec6-vc-dim-f2}. By induction hypothesis, we have
\begin{align*}
  |\mathcal{F}_2| \leq \Phi_{d - 1} (n - 1).
\end{align*}
So we conclude that
\begin{align*}
  |\mathcal{F} (x_1, \ldots, x_n)| = |\mathcal{F}_1 | + |\mathcal{F}_2 | \leq \Phi_d (n - 1) + \Phi_{d - 1} (n - 1) = \Phi_d (n),
\end{align*}
which completes the induction step. The last inequality in the theorem follows from the standard bound on binomial coefficients.



\paragraph{Application of Sauer's lemma: VC dimension of composite classes} Consider two classes of binary functions $\mathcal{F}_1$ and $\mathcal{F}_2$, with VC dimensions $d_1$ and $d_2$ respectively. For any binary operation $\clubsuit$ (e.g. AND, OR, XOR), we can define a composite class of functions
\begin{align*}
  \mathcal{F} = \{f_1 \clubsuit f_2: f_1 \in \mathcal{F}_1, f_2 \in \mathcal{F}_2\}.
\end{align*}
Can we bound the VC dimension of $\mathcal{F}$ in terms of $d_1$ and $d_2$?

Given $(x_1, \ldots, x_n)$, we can consider the set of achievable labelings by $\mathcal{F}$:
\begin{align*}
  \mathcal{F} (x_1, \ldots, x_n) = \{(f_1 (x_1) \clubsuit f_2 (x_1), \ldots, f_1 (x_n) \clubsuit f_2 (x_n)): f_1 \in \mathcal{F}_1, f_2 \in \mathcal{F}_2\}.
\end{align*}
The cardinality of this set is bounded by
\begin{align*}
  |\mathcal{F} (x_1, \ldots, x_n)| \leq |\mathcal{F}_1 (x_1, \ldots, x_n)| \cdot |\mathcal{F}_2 (x_1, \ldots, x_n)| \leq \Big( \frac{e n}{d_1} \Big)^{d_1} \cdot \Big( \frac{e n}{d_2} \Big)^{d_2}.
\end{align*}
On the other hand, let $d$ be the VC dimension of $\mathcal{F}$. By definition, we have
\begin{align*}
  2^d = \max_{x_1, \ldots, x_d} |\mathcal{F} (x_1, \ldots, x_d)| \leq \Big( \frac{e d}{d_1} \Big)^{d_1} \cdot \Big( \frac{e d}{d_2} \Big)^{d_2}.
\end{align*}
Solving for $d$ through this inequality, there exists a universal constant $c > 0$ such that
\begin{align*}
  d \leq c (d_1 + d_2).
\end{align*}


\paragraph{Additional examples of VC dimension:} Consider the class of linear classifiers in $\mathbb{R}^d$:
\begin{align*}
  \mathcal{F} = \Big\{f_\theta (x) = \bm{1}_{\theta^\top x \geq 0}: \theta \in \mathbb{R}^d\Big\}.
\end{align*}
We claim that $\text{VCdim} (\mathcal{F}) = d$. For the upper bound, consider any $(d + 1)$ points in $\mathbb{R}^d$: $x_1, \ldots, x_{d + 1}$. For any possible labeling $(b_1, b_2, \cdots, b_{d + 1}) \in \{0, 1\}^{d + 1}$, suppose that they can be realized by some linear classifier $f_\theta$.

On the other hand, we note that these $(d + 1)$ points are linearly dependent in $\mathbb{R}^d$. That means there exists some not-all-zero coefficients $(\alpha_1, \ldots, \alpha_{d + 1})$ such that
\begin{align*}
  \sum_{i = 1}^{d + 1} \alpha_i x_i = 0.
\end{align*}
Multiplying all coefficients by $-1$ if necessary, we may assume that $\alpha_j < 0$ for some $j$. Let us consider the labeling $b_i = \bm{1}_{\alpha_i \geq 0}$, and let $\theta$ be the linear classifier that realizes this labeling. Then we have
\begin{align*}
  \sum_{i = 1}^{d + 1} \alpha_i \theta^\top x_i = \theta^\top \sum_{i = 1}^{d + 1} \alpha_i x_i = 0, \quad \text{and} \quad \alpha_i \theta^\top x_i \geq 0, \quad \forall i = 1, \ldots, d + 1.
\end{align*}
The only possibility is that $\alpha_i \theta^\top x_i = 0$ for all $i = 1, \ldots, d + 1$. For each $i$, either $\alpha_i = 0$ or $\theta^\top x_i = 0$.

However, $\alpha_j < 0$ gives $b_j = 0$, so $\theta^\top x_j < 0$ and hence $\alpha_j \theta^\top x_j > 0$, a contradiction. Thus no set of $d + 1$ points can be shattered, and $\text{VCdim} (\mathcal{F}) \leq d$.

For the lower bound, take the standard basis vectors $e_1, \ldots, e_d$. For any labeling $(b_1, \ldots, b_d) \in \{0, 1\}^d$, choose $\theta = (2b_1 - 1, \ldots, 2b_d - 1)^\top$. Then $f_\theta (e_i) = b_i$ for every $i$. Hence these $d$ points are shattered, proving $\text{VCdim} (\mathcal{F}) = d$.

\subsection{VC dimension and covering number}
Recall that the goal is to bound the $\mathbb{L}^2 (Q)$ covering number of a function class $\mathcal{F}$ for any distribution $Q$ (for binary classes, $G \equiv 1$ is an envelope function). We have the following result:
\begin{theorem}
  For any $Q$, we have
  \begin{align*}
    \log N (\mathcal{F}; \mathbb{L}^2 (Q), \varepsilon) \leq c \cdot \text{VCdim} (\mathcal{F}) \cdot \log (1/\varepsilon),
  \end{align*}
  for some universal constant $c > 0$.
\end{theorem}
Proof: Let $f_1, f_2, \cdots, f_N$ be a maximal $\varepsilon$-packing of $\mathcal{F}$ with respect to the $\mathbb{L}^2 (Q)$ metric. For any $i, j \in [N]$, we have
\begin{align*}
  \varepsilon^2 \leq \vecnorm{f_i - f_j}{\mathbb{L}^2 (Q)}^2 = \int |f_i (x) - f_j (x)|^2 d Q(x) = Q (\{x: f_i (x) \neq f_j (x)\}).
\end{align*}
The idea is to use probabilistic method: we draw $n$ i.i.d. samples $X_1, \ldots, X_n \simiid Q$. For any $i, j \in [N]$, we have
\begin{align*}
  \Prob \Big( \big( f_i(X_1), f_i (X_2), \ldots, f_i (X_n) \big) = \big( f_j(X_1), f_j (X_2), \ldots, f_j (X_n) \big) \Big) \leq (1 - \varepsilon^2)^n \leq \exp (- n \varepsilon^2).
\end{align*}
Then we can take union bound over all pairs $(i, j)$, and we have
\begin{align*}
  \Prob \Big( \exists i, j \in [N], \text{ not distinguishable under} \big\{X_1, \ldots, X_n\big\} \Big) \leq \binom{N}{2} \exp (- n \varepsilon^2) \leq N^2 \exp (- n \varepsilon^2).
\end{align*}
We choose $n = \frac{4 \log N}{\varepsilon^2}$, then the right-hand-side is less than $1/2$. That means there exists a set of $n$ points $\{x_1, \ldots, x_n\}$ such that all the functions in the packing set are distinguishable on this set of points. That means
\begin{align*}
  N = \Big| \Big\{ (f_i (x_1), \ldots, f_i (x_n)) : i \in [N]\Big\} \Big| \leq \Big| \Big\{ (f (x_1), \ldots, f (x_n)) : f \in \mathcal{F}\Big\} \Big| \leq \Phi_d (n) \leq \Big( \frac{e n}{d} \Big)^d = \Big( \frac{4 e \log N}{d \varepsilon^2} \Big)^d.
\end{align*}
Solving for $N$, we have
\begin{align*}
  \log N \leq c d \log (1/\varepsilon),
\end{align*}
for some universal constant $c > 0$.

e.g. for linear classifiers in $\real^d$, $\mathcal{F} = \{ f(x) = \bm{1}_{\theta^\top x > 0} : \theta \in \real^d  \}$, we have
\begin{align*}
  \Exs \big[ \sup_{f \in \mathcal{F}}|(P_n - P) f| \big] \leq \frac{c}{\sqrt{n}} \int_0^1 \sqrt{d \log (1/ \varepsilon)} d \varepsilon \leq c' \sqrt{\frac{d}{n}}.
\end{align*}
This can also be used to bound the generalization error (gap between training and test error) of linear classifiers in $\real^d$.

\subsection{VC subgraph dimension for real-valued function classes}
Idea: to extend the notion of VC dimension to real-valued function classes.

\begin{definition}[VC subgraph dimension]
  For a real-valued function class $\mathcal{F}$, we define
  \begin{align*}
    \text{VCdim} (\mathcal{F}) := \text{VCdim} (\{ (t, x) \mapsto \bm{1}_{f(x) \geq t}: f \in \mathcal{F}\}).
  \end{align*}
\end{definition}
Intuitively, we augment one dimension to the input space to transform the real-valued function class into a binary function class, while keeping its complexity.

e.g. consider the class of linear functions in $\real^d$:
\begin{align*}
  \mathcal{F} = \{ f_\theta (x) = \theta^\top x: \theta \in \real^d\}.
\end{align*}
We note that $t_i \geq f_\theta (x_i)$ is equivalent to the fact that
\begin{align*}
  \begin{bmatrix}
    -\theta^\top & 1
  \end{bmatrix} \begin{bmatrix}
    x_i\\
    t_i
  \end{bmatrix} \geq 0,
\end{align*}
which becomes a subclass of the $(d + 1)$-dimensional linear classifiers. So we have $\text{VCdim} (\mathcal{F}) \leq d + 1$.

Furthermore, suppose that $\phi$ is a strictly increasing function, we can consider a generalized linear function class
\begin{align*}
  \mathcal{F} = \{ f_\theta (x) = \phi(\theta^\top x): \theta \in \real^d\}.
\end{align*}
We note that $t_i \geq f_\theta (x_i)$ is equivalent to the fact that
\begin{align*}
  \begin{bmatrix}
    -\theta^\top & 1
  \end{bmatrix}
  \begin{bmatrix}
    x_i\\
    \phi^{-1} (t_i)
  \end{bmatrix} \geq 0,
\end{align*}
which is also a subclass of the $(d + 1)$-dimensional linear classifiers.

The nice thing about VC subgraph dimension is that it bounds the $\mathbb{L}^2 (Q)$ covering number with respect to an envelope function, which is exactly what we need for the general empirical process bound. We have the following result:
\begin{theorem}
  Let $\mathcal{F}$ be a real-valued function class with envelope function $G$. Then we have
  \begin{align*}
   \sup_{Q} \log N (\mathcal{F}; \mathbb{L}^2 (Q), \varepsilon \|G\|_{\mathbb{L}^2 (Q)}) \leq \big( c_1 / \varepsilon \big)^{c_2 \text{VCdim} (\mathcal{F})},
  \end{align*}
  for some universal constants $c_1, c_2 > 0$.
\end{theorem}

Proof: for a pair of functions $f, g \in \mathcal{F}$, we have
\begin{align*}
  \vecnorm{f - g}{\mathbb{L}^2 (Q)}^2 &= \int |f(x) - g(x)|^2 d Q(x) \leq 2 \int G(x) |f(x) - g(x)| d Q(x)\\
  &=2 \int G(x) \int | \bm{1}_{t \leq f(x)} - \bm{1}_{t \leq g(x)} | dt~ d Q(x)\\
  &= 2 \iint_{(t, x): |t| \leq G(x)} G (x) \cdot | \bm{1}_{t \leq f(x)} - \bm{1}_{t \leq g(x)} | dt~ d Q(x).
\end{align*}
(The idea is to relate the $\mathbb{L}^2 (Q)$ distance between $f$ and $g$ to the the indicator functions)

By Cauchy--Schwarz inequality, we have
\begin{align*}
  &2 \iint_{(t, x): |t| \leq G(x)} G (x) \cdot | \bm{1}_{t \leq f(x)} - \bm{1}_{t \leq g(x)} | dt~ d Q(x)\\
  &\leq 2 \sqrt{\iint_{(t, x): |t| \leq G(x)} \big( \sqrt{G (x)} \big)^2 dt~ d Q(x)} \cdot \sqrt{\iint_{(t, x): |t| \leq G(x)} \big( \sqrt{G(x)}  | \bm{1}_{t \leq f(x)} - \bm{1}_{t \leq g(x)} | \big)^2 dt~ d Q(x)}.
\end{align*}
For the first integral, we note that
\begin{align*}
  \iint_{(t, x): |t| \leq G(x)} \big( \sqrt{G (x)} \big)^2 dt~ d Q(x) = 2 \int G^2 (x) d Q(x) = 2 \|G\|_{\mathbb{L}^2 (Q)}^2.
\end{align*}
For the second integral, it is
\begin{align*}
  \iint_{(t, x): |t| \leq G(x)}  | \bm{1}_{t \leq f(x)} - \bm{1}_{t \leq g(x)} |^2  G(x) dt~ d Q(x),
\end{align*}
which is the $\mathbb{L}^2 (G (x) dt d Q(x))$-distance of two indicator functions. Indeed, we can construct a probability measure
\begin{align*}
  d \widetilde{Q} (t, x) \propto G(x)dt~ d Q(x).
\end{align*}
We have
\begin{align*}
  \iint_{(t, x): |t| \leq G(x)}  | \bm{1}_{t \leq f(x)} - \bm{1}_{t \leq g(x)} |^2  G(x) dt~ d Q(x) = \vecnorm{\bm{1}_{t \leq f(x)} - \bm{1}_{t \leq g(x)}}{\mathbb{L}^2 (\widetilde{Q})}^2 \cdot 2 \|G\|_{\mathbb{L}^2 (Q)}^2.
\end{align*}
That means an $\varepsilon$-covering of the class of indicator functions $\{(t, x) \mapsto \bm{1}_{t \leq f(x)}: f \in \mathcal{F}\}$ with respect to the $\mathbb{L}^2 (\widetilde{Q})$ metric leads to a $2 \vecnorm{G}{\mathbb{L}^2 (Q)} \sqrt{\varepsilon}$-covering of the class $\mathcal{F}$ with respect to the $\mathbb{L}^2 (Q)$ metric. By the previous result on VC dimension, we have the final bound.



\end{document}
